\chapter{Discusión y conclusiones}
\label{cap:conclusiones}

\section{Discusión}

La discusión sigue el recorrido del trabajo: el fenómeno, los atributos, la optimización y la evaluación. En cada punto
se contrasta lo que predecía la teoría con lo que mostraron los datos.

\subsection{El fenómeno y los atributos de Fourier}

El modelo del Capítulo~\ref{cap:fenomeno} describe la crisis como una descarga rítmica, con fundamental $f_0$ y
armónicos (ecuación~\eqref{eq:descarga}), que se suma a la actividad de fondo y queda multiplicada por una ganancia
propia de cada paciente y canal (ecuación~\eqref{eq:modelo_eeg}). Los datos lo confirmaron en sus tres aspectos:
\begin{itemize}
  \item \textbf{Aumento de potencia concentrado en una banda.} En Bonn el set ictal tiene de 10 a 100 veces más potencia
        entre 3 y 25~Hz, con máximo cerca de 6~Hz, y la banda dominante pasa a ser theta (Tabla~\ref{tab:bonn_bandas}).
        Con solo 20 atributos de Welch, la regresión logística llega a 97,4~\% de exactitud en tres clases y el SVM con
        núcleo RBF, a 98,4~\% (Tabla~\ref{tab:bonn_cv}).
  \item \textbf{La banda cambia entre pacientes.} En CHB-MIT la banda con mayor aumento es delta en chb01 y chb08, y alfa y
        beta en chb05 y chb10 (Tabla~\ref{tab:dif_log}). No hay una ``frecuencia de la crisis'' universal. Esto motiva usar
        varias bandas como atributos y, en la red, convoluciones que se desplazan a lo largo de la frecuencia.
  \item \textbf{La ganancia de cada paciente es un término aditivo en escala logarítmica.} Restar la mediana del paciente lo
        elimina (ecuación~\eqref{eq:calibracion}). Con la partición fija, el AUC de la regresión logística pasa de 0,80 a
        0,93 y las falsas alarmas, de 3,3 a 1,2 por hora (Tabla~\ref{tab:fija_convexos}). Es la mejora más grande del
        trabajo y la produce una operación de una línea, deducida del modelo.
\end{itemize}

El modelo también tiene límites. Las tres crisis de chb03 que la red no detecta tienen un aumento de potencia por banda
\emph{mayor} que las que sí detecta. Un detector que solo mira cuánta potencia hay las encuentra; uno que aprendió la
\emph{forma} espectro-temporal de 13 crisis de otros tres pacientes, no. La ecuación~\eqref{eq:descarga} describe una
descarga estacionaria, y las crisis reales evolucionan en frecuencia y en distribución espacial a lo largo de su duración
(Figura~\ref{fig:chb01_espectrograma}).

\subsection{Optimalidad local y global}

La teoría del Capítulo~\ref{cap:optimizacion} se verificó numéricamente en los dos regímenes:
\begin{itemize}
  \item \textbf{Problema convexo.} La regresión logística regularizada es fuertemente convexa
        (Proposición~\ref{prop:logistica_convexa}). Desde 20 puntos iniciales a distancia media 13 del óptimo, el descenso por
        gradiente con paso $1/L$ llega al mismo valor $J^* = 0{,}07699786$ con distancia máxima $3\times10^{-7}$ entre
        soluciones. La convergencia es lineal, con una pendiente compatible con $\kappa \approx 528$
        (Teorema~\ref{teo:convergencia}). Todo mínimo local es global (Teorema~\ref{teo:local_global}) y una sola corrida
        alcanza.
  \item \textbf{Problema no convexo.} Un perceptrón de 2 neuronas termina en 8 valores distintos de la pérdida según la
        semilla, con una diferencia de un factor 20. Con 4 neuronas todas las semillas llegan a la misma pérdida, pero con
        pesos distintos, y en el punto medio entre dos soluciones la pérdida es 14 veces mayor: una prueba directa de no
        convexidad.
\end{itemize}

La consecuencia práctica se vio en la CNN. Con el mismo optimizador y la misma tasa, el AUC de prueba de AdamW va de 0,873
a 0,960 según la semilla (Tabla~\ref{tab:comparativa}). Comparar optimizadores con una corrida de cada uno habría dado
cualquier orden. Por eso se usaron 5 semillas y una prueba no paramétrica, y aun así la potencia es limitada: con $n = 5$
por grupo, el menor valor $p$ posible es 0,008.

\subsection{Lagrange, KKT y simplex}

Los multiplicadores de Lagrange aparecieron en tres lugares:
\begin{itemize}
  \item En el \textbf{dual del SVM}, donde las condiciones KKT se verificaron en la solución numérica (Tabla~\ref{tab:kkt_verif}).
        Los multiplicadores clasifican los puntos en los tres casos de la holgura complementaria (Tabla~\ref{tab:svm_kkt}),
        y la brecha de dualidad es de $1{,}6\times10^{-7}$.
  \item En la \textbf{softmax}, cuyas salidas viven en el simplex de probabilidad (Figura~\ref{fig:bonn_simplex}): la
        restricción $\sum_k p_k = 1$ es la que hace que el gradiente sea $\mathbf{p} - \mathbf{y}$ (ecuación~\eqref{eq:grad_ce}).
  \item En el \textbf{umbral de Neyman-Pearson}, donde el multiplicador $\lambda$ de la restricción de falsos positivos es la
        pendiente de la curva ROC en el punto de operación (ecuación~\eqref{eq:pendiente_roc}) y tiene una interpretación
        concreta: cuántas crisis se ganan por cada falsa alarma adicional. En el barrido de $\alpha$ de la CNN, esa pendiente
        va de 2 crisis por cada 0,19 falsas alarmas por hora a 1 crisis por cada 7 (Tabla~\ref{tab:barrido_alfa}).
\end{itemize}
La versión discreta de Neyman-Pearson es un programa lineal (ecuación~\eqref{eq:np_pl}), y el Teorema~\ref{teo:vertice}
explica por qué su solución es un umbral determinístico salvo, a lo sumo, una ventana aleatorizada: el óptimo está en un
vértice del poliedro factible, que es justamente lo que recorre el simplex.

\subsection{La comparativa de optimizadores}

El resultado principal de la comparativa no es cuál optimizador gana, sino \emph{por qué}. Con la tasa elegida por
validación, SGD supera significativamente a Momentum, Nesterov, Adam y AdamW. Esto parecía confirmar la observación de
\textcite{wilson2017} de que los métodos adaptativos generalizan peor. El experimento exploratorio con la tasa más alta de
cada grilla muestra otra cosa: SGD, RMSprop, Adam y AdamW quedan entre 0,970 y 0,976, sin diferencias significativas
(Tabla~\ref{tab:tasa_alta}). La ventaja de SGD era, sobre todo, una ventaja de la tasa de aprendizaje.

Hay dos lecturas posibles y no son excluyentes:
\begin{itemize}
  \item \textbf{La tasa importa más que el algoritmo.} Con el paso adecuado, los seis optimizadores llegan a soluciones que
        generalizan parecido. El análisis de la cuadrática (Proposición~\ref{prop:momentum}) muestra que el momento acelera
        y estabiliza, pero no cambia \emph{a qué} mínimo se llega, y en un problema con muchos mínimos de pérdida parecida
        es esto último lo que determina la generalización.
  \item \textbf{Los pasos grandes favorecen soluciones más robustas.} Todos los optimizadores mejoraron con la tasa alta.
        Es consistente con la idea de que el ruido de los pasos grandes evita los mínimos angostos, aunque este trabajo no
        midió la curvatura de las soluciones y no puede confirmarlo.
\end{itemize}

Hubo además un efecto de diseño. El paciente de validación quedó saturado (AUC entre 0,996 y 0,999 para las 42 redes) y su
AUC no predijo el de prueba ($\rho = -0{,}30$). El criterio de elección no tenía información, y la tasa elegida fue en buena
parte azar. Esa fue la causa de que se compararan pares (optimizador, tasa) que no eran equivalentes.

\subsection{Qué métrica usar}

La comparación final entre la CNN y la regresión logística calibrada resume la lección del Capítulo~\ref{cap:metricas}.
La logística tiene mayor AUC (0,934 contra 0,872), pero la red detecta 11 de 14 crisis contra 7, con falsas alarmas
parecidas y menos de la mitad de latencia. Con una tasa de falsos positivos de $10^{-3}$, la red detecta el 54~\% de las
ventanas ictales y la logística, el 2~\%.

El AUC promedia toda la curva ROC con el mismo peso, y la mayor parte de esa curva corresponde a tasas de falsas alarmas
(decenas o cientos por hora) que no tienen uso clínico. Con 900 ventanas por hora, la zona útil es
$\text{FPR} \lesssim 10^{-3}$. Ahí se decide la calidad de un detector, y ahí conviene medir: sensibilidad con una tasa
fija de falsos positivos, AUC parcial, o directamente la curva de sensibilidad por evento contra falsas alarmas por hora
(Figura~\ref{fig:sens_fa}). La exactitud es todavía peor: un detector que nunca avisa acierta más del 99,6~\% de las ventanas.

\section{Limitaciones}

\begin{enumerate}
  \item \textbf{Bonn no identifica a los pacientes.} Los 100 segmentos de cada set provienen de pocos sujetos, y la validación
        cruzada estratificada pone segmentos del mismo paciente en entrenamiento y en prueba. Las exactitudes de 97--98~\%
        son optimistas y no miden generalización a pacientes nuevos. Además, los segmentos se seleccionaron a mano y están
        libres de artefactos.
  \item \textbf{Pocos pacientes de CHB-MIT.} Se usaron 6 de los 23 pacientes de la base y 34 crisis, todos niños. La
        evaluación final de la CNN descansa en 2 pacientes y 14 crisis, y las diferencias entre optimizadores, en 3 crisis de
        un solo paciente (chb03\_34 a chb03\_36). Los intervalos de confianza de cualquier proporción con 14 eventos son
        anchos: 11/14 es compatible con una sensibilidad real de entre 0,49 y 0,95 (intervalo exacto de Clopper-Pearson del
        95~\%).
  \item \textbf{Un solo paciente de validación.} Para la CNN se usó una partición fija y no validación cruzada por paciente,
        por costo de cómputo. La validación quedó saturada y no sirvió para elegir la tasa ni el modelo final.
  \item \textbf{Análisis \emph{a posteriori}.} El análisis por crisis y el experimento con la tasa alta se diseñaron después de
        mirar la prueba. Sus conclusiones son hipótesis, no resultados confirmados, y habría que validarlas con otros
        pacientes.
  \item \textbf{La calibración usa el futuro.} La mediana de cada paciente de prueba se calcula sobre todo su registro. Un
        sistema en línea tendría que estimarla con las primeras horas o con una ventana móvil. La mediana es robusta y, con
        crisis que ocupan menos del 1~\% del tiempo, el efecto debería ser chico, pero no se midió.
  \item \textbf{El umbral no se transfiere bien entre pacientes.} Neyman-Pearson garantiza la tasa de falsos positivos solo si
        la distribución del puntaje sin crisis es la misma en la calibración y en el uso (Observación~\ref{obs:np_supuesto}).
        Entre pacientes no lo es: en chb03, con $\alpha = 10^{-3}$, la tasa real fue siete veces mayor.
  \item \textbf{Las ventanas de borde se excluyen.} Las ventanas que mezclan crisis y actividad normal no se usan en el
        entrenamiento ni en las métricas por ventana. Las métricas por evento sí cubren el registro completo.
  \item \textbf{Cómputo limitado.} Sin GPU, la red es chica (64\,481 parámetros), se entrena 30 épocas, la grilla de tasas tiene
        3 valores y no se exploraron otros hiperparámetros (tamaño de lote, momento, $\beta_2$, decaimiento, programación de
        la tasa) ni otras arquitecturas.
\end{enumerate}

\section{Trabajo futuro}

\begin{itemize}
  \item \textbf{Validación cruzada por paciente para la CNN}, tanto para elegir la tasa como para estimar el desempeño. Con
        6 pacientes y 92~s por corrida, un LOPO anidado (5 particiones internas, 3 tasas, 6 optimizadores) son 540 corridas,
        unas 14 horas de CPU con una semilla, o menos de una hora con GPU.
  \item \textbf{Más pacientes de CHB-MIT}, y evaluación en otra base (por ejemplo, la de Siena o la del Hospital de la
        Universidad de Temple) para medir generalización entre centros y equipos.
  \item \textbf{Calibración en línea}: mediana móvil o estimada con la primera hora de registro, y un umbral adaptado al
        paciente con un período inicial sin crisis. Atacaría directamente el problema de chb03.
  \item \textbf{Optimizar la métrica que importa}: elegir hiperparámetros por sensibilidad con una tasa fija de falsos
        positivos o por AUC parcial, en lugar del AUC completo.
  \item \textbf{Entender la tasa alta}: medir la curvatura de las soluciones (autovalor máximo de la Hessiana, por iteración de
        la potencia con productos Hessiana-vector) para ver si las redes entrenadas con paso grande están, efectivamente, en
        mínimos más planos.
  \item \textbf{Combinar modelos}: la logística calibrada puntúa alto justamente las tres crisis que la red no ve. Un promedio
        de los puntajes, o usar la potencia por banda como entrada adicional de la red, podría tomar lo mejor de ambos.
  \item \textbf{Tiempo y espacio}: modelos que miren varias ventanas seguidas (recurrentes o convoluciones temporales) y que
        tengan en cuenta la posición de los electrodos.
\end{itemize}

\section{Conclusiones}

\begin{enumerate}
  \item La crisis epiléptica tiene una firma espectral clara: un aumento de potencia de uno a dos órdenes de magnitud en una
        banda que depende del paciente. La transformada de Fourier y sus estimadores (Welch, espectrograma) la convierten en
        atributos con los que un modelo lineal ya separa bien las crisis típicas.
  \item El modelo matemático del EEG no es solo descriptivo: de él se deduce la calibración por paciente, que es la mejora más
        grande del trabajo (AUC de 0,80 a 0,93 con la regresión logística en pacientes nuevos).
  \item La teoría de optimización predice correctamente lo observado. En el problema convexo, todos los puntos iniciales
        llegan al mismo óptimo con la velocidad que da el número de condición. En el no convexo, el resultado depende de la
        semilla, lo que obliga a repetir corridas y a usar pruebas estadísticas para comparar.
  \item Los multiplicadores de Lagrange conectan los tres problemas con restricciones del trabajo: el dual del SVM, la softmax
        y el umbral de Neyman-Pearson. Este último, que es un programa lineal en su versión discreta, es la herramienta
        adecuada para fijar el punto de operación de un detector con clases tan desbalanceadas.
  \item En la comparativa de optimizadores, la tasa de aprendizaje pesó más que el algoritmo: con la tasa adecuada, SGD,
        RMSprop, Adam y AdamW no se distinguen. Para que esa conclusión sea firme falta repetirla con validación cruzada por
        paciente.
  \item La red convolucional es el mejor detector en la zona clínicamente útil: con una falsa alarma cada 5 horas detecta 10 de
        14 crisis, contra 6 de la regresión logística con casi el triple de falsas alarmas. Esto pasa a pesar de que su AUC es
        menor y de que el protocolo eligió una de sus peores corridas.
  \item Para un detector de crisis, la sensibilidad, la especificidad y el AUC por ventana no alcanzan. Lo que hay que informar
        es cuántas crisis se detectan, con cuántas falsas alarmas por hora y con qué demora, sobre pacientes que el modelo no
        vio.
\end{enumerate}
